
@@M0@@
Le
groupe @@M1@@ opère sur @@M2@@, sur @@M3@@ muni des faisceaux @@M4@@ et sur @@M5@@. Si on fait le produit de ces objets avec @@M6@@, sous @@M7@@, on obtient le faisceau @@M8@@ associé à @@M9@@, le schéma @@M10@@ des drapeaux de @@M11@@, et le schéma @@M12@@. On notera les faisceaux canoniques sur @@M13@@ par les mêmes lettres que leurs homologues sur @@M14@@. On a donc le diagramme :
@@M15@@
On remarque que @@M16@@ est un torseur de base @@M17@@, de groupe structural @@M18@@, et que @@M19@@ est un torseur de base @@M20@@, de groupe structural @@M21@@, défini par la @@M22@@-Algèbre :
@@M23@@

L'homomorphisme @@M24@@ se factorise donc en :
@@M25@@
D'après 3.1, on a @@M26@@. Soit @@M27@@ la classe de @@M28@@ dans @@M29@@. Alors d'après 2.2, @@M30@@ est le quotient de @@M31@@ par le sous-module engendré par les @@M32@@. Si @@M33@@ est l'idéal de @@M34@@ engendré par les @@M35@@, on a donc :
@@M36@@
